% Luc Destombe --- licence CC BY-NC-SA 4.0
% https://creativecommons.org/licenses/by-nc-sa/4.0/deed.fr
% Fichier autonome : compiler avec pdflatex, deux fois (signets, nombre de pages).
\def\lyceeversion{1.0}
\documentclass[a4paper,10pt]{article}
\makeatletter
\newif\iflycee@niveaux
\newif\iflycee@cm
\newif\iflycee@palatino

\RequirePackage[utf8]{inputenc}
\RequirePackage[T1]{fontenc}
\RequirePackage[french]{babel}
\RequirePackage{etoolbox}

\newcommand{\ShorthandsOff}{\@ifundefined{shorthandoff}{}{\shorthandoff{;:!?}}}
\newcommand{\ShorthandsOn}{\@ifundefined{shorthandon}{}{\shorthandon{;:!?}}}
\ShorthandsOff
\AtBeginDocument{\ShorthandsOn}
\AtBeginEnvironment{tikzpicture}{\ShorthandsOff}

\RequirePackage{geometry}
\RequirePackage{fancyhdr}
\RequirePackage{lastpage}
\RequirePackage{multicol}
\RequirePackage{array}
\RequirePackage{enumitem}
\RequirePackage{xcolor}
\RequirePackage{needspace}

\RequirePackage{amsmath,amssymb}
\iflycee@cm\else\RequirePackage{fourier}\fi
\DeclareMathAlphabet{\mathcal}{OMS}{cmsy}{m}{n}
\renewcommand{\ttdefault}{lmtt}
\begingroup\catcode`\_=\active \catcode`\^=\active
\global\let_\sb \global\let^\sp
\endgroup
\AtBeginDocument{\catcode`\_=12 \mathcode`\_="8000
  \catcode`\^=12 \mathcode`\^="8000 }
\newlength{\retraitcalculs}
\setlength{\retraitcalculs}{2em}

\RequirePackage{tikz}
\usetikzlibrary{arrows.meta,calc,babel,shapes.misc}
\RequirePackage[most]{tcolorbox}
\tcbuselibrary{skins,breakable}

\definecolor{coulDef}{HTML}{1F4E79}
\definecolor{coulProp}{HTML}{7030A0}
\definecolor{coulRem}{HTML}{C55A11}
\definecolor{coulEx}{HTML}{2E7D32}
\definecolor{coulExo}{HTML}{B03A2E}
\definecolor{coulPart}{HTML}{1F4E79}
\definecolor{coulMeth}{HTML}{1B6E4F}
\definecolor{coulCourbe}{HTML}{0B5ED7}
\definecolor{coulCourbeB}{HTML}{C00000}
\definecolor{coulCourbeC}{HTML}{2E7D32}
\definecolor{coulGrille}{HTML}{8C8C8C}
\definecolor{coulRep}{HTML}{1B6E4F}
\definecolor{coulBar}{HTML}{2E7D32}

\newcommand{\R}{\mathbb{R}}
\newcommand{\Rs}{\mathbb{R}^{*}}
\newcommand{\N}{\mathbb{N}}
\newcommand{\Ns}{\mathbb{N}^{*}}
\newcommand{\Z}{\mathbb{Z}}
\newcommand{\Q}{\mathbb{Q}}
\renewcommand{\le}{\leqslant}
\renewcommand{\ge}{\geqslant}
\renewcommand{\approx}{\simeq}
\newcommand{\vect}[1]{\overrightarrow{#1}}
\newcommand{\norme}[1]{\left\lVert #1 \right\rVert}
\newcommand{\intervalle}[2]{\left[#1\,;#2\right]}
\RequirePackage{cancel}
\renewcommand{\CancelColor}{\color{coulExo}}
\newcommand{\fois}[1]{\times{\color{coulExo}#1}}
\newcommand{\barre}[1]{\cancel{#1}}
\newlength{\lycee@eqh}\newlength{\lycee@eqd}
\newsavebox{\lycee@eqa}\newsavebox{\lycee@eqb}\newsavebox{\lycee@eqc}
\newcommand{\lycee@eqcol}[1]{\hbox{$\renewcommand{\arraystretch}{1.35}%
  \begin{array}[t]{@{}>{\displaystyle}r@{}>{\displaystyle{}}l@{}}#1\end{array}$}}
\newcommand{\lycee@eqmesure}[1]{%
  \ifdim\ht#1>\lycee@eqh\setlength{\lycee@eqh}{\ht#1}\fi
  \ifdim\dp#1>\lycee@eqd\setlength{\lycee@eqd}{\dp#1}\fi}
\newcommand{\lycee@eqtrait}{\hspace{1.2em}\tikz[baseline=0pt]\draw[dash pattern=on 2pt off 2pt,
  line width=0.5pt] (0,\lycee@eqh) -- (0,-\lycee@eqd);\hspace{1.2em}}
\newcommand{\eqproduit}[2]{\par\smallskip\noindent C'est une équation produit :
  $A\times B=0 \iff A=0$ ou $B=0$.\par\nopagebreak\smallskip
  \begingroup\setlength{\lycee@eqh}{0pt}\setlength{\lycee@eqd}{0pt}%
  \setbox\lycee@eqa=\lycee@eqcol{#1}\setbox\lycee@eqb=\lycee@eqcol{#2}%
  \lycee@eqmesure\lycee@eqa\lycee@eqmesure\lycee@eqb%
  \noindent\hspace*{\retraitcalculs}\usebox\lycee@eqa\lycee@eqtrait\usebox\lycee@eqb\endgroup\par\smallskip}
\newcommand{\eqproduittrois}[3]{\par\smallskip\noindent C'est une équation produit :
  $A\times B\times C=0 \iff A=0$ ou $B=0$ ou $C=0$.\par\nopagebreak\smallskip
  \begingroup\setlength{\lycee@eqh}{0pt}\setlength{\lycee@eqd}{0pt}%
  \setbox\lycee@eqa=\lycee@eqcol{#1}\setbox\lycee@eqb=\lycee@eqcol{#2}\setbox\lycee@eqc=\lycee@eqcol{#3}%
  \lycee@eqmesure\lycee@eqa\lycee@eqmesure\lycee@eqb\lycee@eqmesure\lycee@eqc%
  \noindent\hspace*{\retraitcalculs}\usebox\lycee@eqa\lycee@eqtrait\usebox\lycee@eqb\lycee@eqtrait\usebox\lycee@eqc\endgroup
  \par\smallskip}

\tikzset{grille/.style={coulGrille,line width=0.4pt}}
\newcommand{\quadrillage}[4]{\draw[grille,step=1] (#1,#2) grid (#3,#4);}
\tikzset{
  croix/.style={cross out,draw,line width=0.6pt,inner sep=0pt,minimum size=4pt},
  croixdroite/.style={croix,rotate=45,line width=0.8pt},
}
\newcommand{\pt}[3]{\path (#1) node[croix]{} node[#2,font=\small,inner sep=2.5pt]{$#3$};}

\newcommand{\lycee@repere}[4]{%
  \draw[grille,step=1] (#1,#2) grid (#3,#4);
  \draw[-{Stealth[length=2mm]},thick] (#1,0)--({#3+0.4},0) node[below left=-1pt]{$x$};
  \draw[-{Stealth[length=2mm]},thick] (0,#2)--(0,{#4+0.4}) node[below left=-1pt]{$y$};
  \node[below left,font=\scriptsize,inner sep=1.5pt] at (0,0) {$0$};}
\newcommand{\lycee@gradx}[1]{\foreach \k/\l in {#1}{%
  \draw (\k,0.07)--(\k,-0.07) node[below,font=\scriptsize,inner sep=1.5pt]{$\l$};}}
\newcommand{\lycee@grady}[1]{\foreach \k/\l in {#1}{%
  \draw (0.07,\k)--(-0.07,\k) node[left,font=\scriptsize,inner sep=1.5pt]{$\l$};}}
\newcommand{\lycee@intff}[2]{\left[#1\,;#2\right]}
\newcommand{\lycee@intfo}[2]{\left[#1\,;#2\right[}
\newcommand{\lycee@intof}[2]{\left]#1\,;#2\right]}
\newcommand{\lycee@intoo}[2]{\left]#1\,;#2\right[}
\newcommand{\lycee@ens}[1]{\left\{#1\right\}}
\AtBeginDocument{%
  \providecommand{\repere}{\lycee@repere}%
  \providecommand{\gradx}{\lycee@gradx}%
  \providecommand{\grady}{\lycee@grady}%
  \providecommand{\Cf}{\mathcal{C}\sb{f}}%
  \providecommand{\Cg}{\mathcal{C}\sb{g}}%
  \providecommand{\intff}{\lycee@intff}%
  \providecommand{\intfo}{\lycee@intfo}%
  \providecommand{\intof}{\lycee@intof}%
  \providecommand{\intoo}{\lycee@intoo}%
  \providecommand{\ens}{\lycee@ens}}

\newlist{questions}{enumerate}{3}
\setlist[questions,1]{label=\arabic*\textdegree),leftmargin=*,itemsep=2pt,topsep=3pt}
\setlist[questions,2]{label=\alph*),leftmargin=*,itemsep=1pt,topsep=2pt}
\setlist[questions,3]{label=\roman*),leftmargin=*,itemsep=1pt,topsep=1pt}
\setlist[itemize]{label=\textbullet,leftmargin=*,itemsep=2pt,topsep=3pt}

\newcommand{\besoinplace}[1]{\par\penalty\@M\begingroup
  \dimen@=#1\relax\dimen@ii=\pagegoal\advance\dimen@ii-\pagetotal
  \ifdim\dimen@>\dimen@ii\ifdim\dimen@ii>\z@\vfil\fi\break\fi\endgroup}

\newcommand{\lycee@pied}{\thepage/\pageref{LastPage}}
\newcommand{\entete}[2]{%
  \pagestyle{fancy}\fancyhf{}%
  \fancyhead[L]{\footnotesize\color{coulPart}\textbf{#1}}%
  \fancyhead[R]{\footnotesize\color{coulPart}\textbf{#2}}%
  \fancyfoot[C]{\footnotesize\lycee@pied}%
  \renewcommand{\headrulewidth}{0.4pt}%
  \renewcommand{\headrule}{\hbox to\headwidth{%
    \color{coulPart!40}\leaders\hrule height \headrulewidth\hfill}}}

\RequirePackage{ccicons}
\newcommand{\lycee@licence}{{\scriptsize\color{gray}%
  \copyright\ Luc Destombe --- \ccbyncsa\ CC BY-NC-SA 4.0}}
\newif\iflycee@licence \lycee@licencetrue
\newcommand{\sanslicence}{\lycee@licencefalse}
\AtBeginDocument{\iflycee@licence\AddToHookNext{shipout/foreground}{%
  \put(0,0){\rlap{\kern\dimexpr1in+\hoffset+\oddsidemargin+\textwidth\relax
    \llap{\raisebox{-\dimexpr1in+\voffset+\topmargin+\headheight+\headsep
      +\textheight+\footskip\relax}{\lycee@licence}}}}}\fi}

\newcommand{\lycee@site}{faire-des-maths.fr}
\newcommand{\lycee@version}{{\scriptsize\color{gray}\lycee@site
  \ifdefined\lyceeversion\ --- Version~\lyceeversion\fi}}
\AtBeginDocument{\iflycee@licence\AddToHookNext{shipout/foreground}{%
  \put(0,0){\rlap{\kern\dimexpr1in+\hoffset+\oddsidemargin\relax
    \raisebox{-\dimexpr1in+\voffset+\topmargin+\headheight+\headsep
      +\textheight+\footskip\relax}{\lycee@version}}}}\fi}

\newcommand{\formulesserrees}{%
  \setlength{\abovedisplayskip}{3pt}\setlength{\belowdisplayskip}{3pt}%
  \setlength{\abovedisplayshortskip}{2pt}\setlength{\belowdisplayshortskip}{3pt}}

\setlength{\parindent}{0pt}

\RequirePackage[implicit=false,hidelinks,pdfencoding=auto,psdextra,bookmarksopen,
  bookmarksopenlevel=1,pdfcreator={},pdfproducer={}]{hyperref}
\RequirePackage{bookmark}
\hypersetup{pdfauthor={Luc Destombe},pdfkeywords={CC BY-NC-SA 4.0}}
\newcounter{lycee@signet}
\newcommand{\signet}[3][]{\stepcounter{lycee@signet}%
  \edef\lycee@dest{\if\relax\detokenize{#1}\relax
    signet.\arabic{lycee@signet}\else#1\fi}%
  \hypertarget{\lycee@dest}{}\bookmark[level=#2,dest=\lycee@dest]{#3}}
\pdfstringdefDisableCommands{%
  \def\R{R}\def\N{N}\def\Z{Z}\def\Q{Q}\def\vect#1{#1}\def\Cf{Cf}\def\Cg{Cg}\def\fois#1{×#1}%
  \def\barre#1{#1}\def\intervalle#1#2{[#1;#2]}\def\intff#1#2{[#1;#2]}%
  \def\textsuperscript#1{#1}\def\dfrac#1#2{#1/#2}\def\frac#1#2{#1/#2}%
  \def\vec#1{#1⃗}}%
\RequirePackage[official]{eurosym}

\geometry{top=0.8cm,bottom=1.3cm,left=1cm,right=1cm,footskip=0.6cm}

\pagestyle{fancy}
\fancyhf{}
\renewcommand{\headrulewidth}{0pt}
\fancyfoot[C]{\footnotesize Page \thepage{} / \pageref{LastPage}}

\newcommand{\titrefiche}[2]{%
  \begin{center}
    \Large\textbf{#1}\\[2pt]
    \large\textbf{#2}\\[2pt]
  \end{center}
  \vspace{0.10cm}}

\setlength{\columnseprule}{0.4pt}
\setlength{\columnsep}{15pt}
\renewcommand{\columnseprulecolor}{\color{black!45}}
\setlength{\multicolsep}{2pt}

\newcounter{partiefiche}
\renewcommand{\thepartiefiche}{\Roman{partiefiche}}
\newif\ifapresbandeau
\newcommand{\lycee@bandeau}[1]{%
  \par\penalty-600
  \vspace{0.08cm}%
  \noindent\colorbox{coulPart}{%
    \parbox{\dimexpr\linewidth-2\fboxsep\relax}{%
      \color{white}\bfseries\raggedright\strut#1\strut}}%
  \nopagebreak\par\nopagebreak
  \vspace{0.06cm}\nopagebreak
  \apresbandeautrue}
\newcommand{\partiefiche}[1]{%
  \refstepcounter{partiefiche}%
  \lycee@bandeau{\signet{1}{\thepartiefiche{} --- #1}\thepartiefiche{} --- #1}}
\newcommand{\partiefichesansnum}[1]{\lycee@bandeau{\signet{1}{#1}#1}}

\newcounter{exo}
\newlength{\espaceexo}\setlength{\espaceexo}{7pt}
\newlength{\espacetitre}\setlength{\espacetitre}{2pt}
\newcommand{\exo}[1][\relax]{%
  \par
  \ifapresbandeau\apresbandeaufalse\else\penalty-150\addvspace{\espaceexo}\fi
  \refstepcounter{exo}%
  \noindent\signet[exercice-\arabic{exo}]{2}{Exercice \arabic{exo}}%
  \textbf{\color{coulExo}\underline{Exercice \theexo}}%
  \ifx\relax#1\relax\else\textbf{ : #1}\fi
  \par\nopagebreak\vspace{\espacetitre}\nopagebreak
  \ignorespaces}

\setlist[enumerate]{nosep,leftmargin=*,itemsep=2pt,topsep=2pt}
\setlist[itemize]{nosep,leftmargin=*,topsep=2pt}
\newcommand{\choix}[3]{\par\hspace*{1em}%
  \makebox[0.3\linewidth][l]{a) #1}\makebox[0.3\linewidth][l]{b) #2}c) #3}
\newcommand{\deux}[2]{\par\hspace*{0.6em}%
  \makebox[0.48\linewidth][l]{#1}#2}
\newcommand{\trois}[3]{\par\hspace*{0.6em}%
  \makebox[0.32\linewidth][l]{#1}\makebox[0.32\linewidth][l]{#2}#3}

\renewenvironment{center}{\par\addvspace{3pt}\centering}{\par\addvspace{3pt}}
\formulesserrees
\AtBeginDocument{\formulesserrees}
\clubpenalty=10000
\widowpenalty=10000
\displaywidowpenalty=10000
\brokenpenalty=10000
\setlength{\parskip}{0pt}

\RequirePackage{ragged2e}
\setlength{\RaggedRightRightskip}{0pt plus 1fil}
\AtBeginDocument{\RaggedRight\binoppenalty=10000 \relpenalty=10000 }
\g@addto@macro\@parboxrestore{\RaggedRight}
\makeatother
\geometry{bottom=1cm}

\setlength{\espaceexo}{4pt}
\setlist[enumerate]{nosep,leftmargin=*,itemsep=1pt,topsep=1pt}

\newcommand{\sys}[1]{\left\{\begin{aligned}#1\end{aligned}\right.}
\newcommand{\coord}[2]{\begin{pmatrix}#1\\#2\end{pmatrix}}

\begin{document}

\titrefiche{Seconde --- Systèmes d'équations}{Fiche d'exercices du chapitre}

{\small\itshape Le repère est orthonormé dans toute la fiche.\par}
\vspace{0.12cm}

\begin{multicols}{2}\raggedcolumns

\partiefiche{Systèmes par substitution}

\exo[Être solution d'un système]
On considère le système ${\sys{3x+y &= 11\\ x-2y &= -1}}$\,.

Pour chaque couple $(3\,;2)$, $(4\,;-1)$ et $(1\,;1)$, dire par une phrase s'il est
solution du système, en justifiant.

\exo[Une inconnue déjà isolée]
Résoudre chaque système par substitution.
\begin{center}
  a)~${\sys{y &= x+3\\ 2x+y &= 12}}$ \qquad
  b)~${\sys{x &= 2y-1\\ 3x-4y &= 1}}$\\[4pt]
  c)~${\sys{y &= -2x+5\\ 4x-y &= 7}}$
\end{center}

\exo[Isoler, puis remplacer]
Résoudre chaque système par substitution.
\begin{center}
  a)~${\sys{x+2y &= 5\\ 3x+5y &= 13}}$ \qquad
  b)~${\sys{2x+y &= 7\\ 5x+3y &= 18}}$\\[4pt]
  c)~${\sys{3x-y &= 4\\ 5x-2y &= 6}}$ \qquad
  d)~${\sys{-x+4y &= 6\\ 2x-3y &= 3}}$
\end{center}

\exo[Des solutions en fractions]
Résoudre chaque système par substitution.
\begin{center}
  a)~${\sys{x+3y &= 2\\ 2x-3y &= 1}}$ \qquad
  b)~${\sys{y &= x+1\\ 5x-2y &= 2}}$
\end{center}

\partiefiche{Systèmes par combinaison}

\exo[Des coefficients déjà opposés]
Résoudre chaque système par combinaison.
\begin{center}
  a)~${\sys{2x+3y &= 12\\ 4x-3y &= 6}}$ \qquad
  b)~${\sys{-5x+2y &= -1\\ 5x+4y &= 13}}$\\[4pt]
  c)~${\sys{3x+4y &= 1\\ -3x+2y &= 5}}$
\end{center}

\exo[Multiplier une équation]
Résoudre chaque système par combinaison, en multipliant une seule des deux équations.
\begin{center}
  a)~${\sys{3x+2y &= 8\\ 6x+5y &= 17}}$ \qquad
  b)~${\sys{4x+3y &= 5\\ 2x+5y &= 13}}$\\[4pt]
  c)~${\sys{5x+2y &= 1\\ 3x-4y &= 11}}$
\end{center}

\exo[Multiplier les deux équations]
Résoudre chaque système par combinaison.
\begin{center}
  a)~${\sys{3x+5y &= -1\\ 5x+3y &= 9}}$ \qquad
  b)~${\sys{2x+3y &= 4\\ 5x+4y &= 3}}$\\[4pt]
  c)~${\sys{4x-3y &= 13\\ 3x+2y &= 14}}$
\end{center}

\exo[Choisir sa méthode]
Pour chaque système, dire par une phrase quelle méthode est la plus rapide et pourquoi,
puis le résoudre.
\begin{center}
  a)~${\sys{-x+y &= -7\\ 2x+y &= 8}}$ \qquad
  b)~${\sys{x &= 3y+1\\ 2x-5y &= 4}}$\\[4pt]
  c)~${\sys{4x+3y &= 10\\ 6x-5y &= -4}}$ \qquad
  d)~${\sys{2x+3y &= 2\\ 4x-3y &= 1}}$
\end{center}

\partiefiche{Problèmes et systèmes}

\exo[Mangas et bandes dessinées]
Lou achète un manga et deux bandes dessinées pour $38$~€. Sam achète quatre mangas et
trois bandes dessinées pour $77$~€. Tous les mangas ont le même prix, toutes les bandes
dessinées aussi. On note $x$ le prix d'un manga et $y$ celui d'une bande dessinée, en
euros.
\begin{enumerate}
  \item Écrire l'équation qui traduit l'achat de Lou, puis celle qui traduit l'achat
        de Sam.
  \item Résoudre le système formé par ces deux équations.
  \item Donner le prix d'un manga et celui d'une bande dessinée.
\end{enumerate}

\exo[Imprimante 3D]
Un club de technologie imprime en 3D des porte-clés, qui demandent chacun $20$~g de
filament, et des figurines, qui en demandent $50$~g. En une semaine, il a imprimé $40$
objets et utilisé $1\,220$~g de filament.

Combien de porte-clés et combien de figurines le club a-t-il imprimés ?

\exo[Sortie scolaire]
Pour une sortie scolaire, deux autocaristes proposent leurs tarifs.
\begin{center}
\renewcommand{\arraystretch}{1.3}
\begin{tabular}{|l|c|c|c|}
\hline
 & par élève & par adulte & total \\ \hline
Autocariste A & $12$~€ & $20$~€ & $680$~€ \\ \hline
Autocariste B & $10$~€ & $30$~€ & $620$~€ \\ \hline
\end{tabular}
\end{center}
Combien d'élèves et combien d'adultes participent à la sortie ?

\partiefiche{Systèmes et intersection de droites}

\exo[Point d'intersection]
Pour chaque couple de droites, démontrer que les droites sont sécantes, puis calculer
les coordonnées de leur point d'intersection.
\begin{enumerate}[label=\alph*)]
  \item $(d_{1})$ : $x+y-5=0$ et $(d_{2})$ : $2x-y-1=0$.
  \item $(d_{3})$ : $y=3x-2$ et $(d_{4})$ : $y=-x+6$.
  \item $(d_{5})$ : $3x+2y-4=0$ et $(d_{6})$ : $x-4y+8=0$.
\end{enumerate}

\exo[Parallèles ou sécantes ?]
Pour chaque couple de droites, dire par une phrase si les droites sont parallèles ou
sécantes ; si elles sont sécantes, calculer les coordonnées de leur point
d'intersection.
\begin{enumerate}[label=\alph*)]
  \item $(d_{1})$ : $2x-6y+1=0$ et $(d_{2})$ : $-x+3y+4=0$.
  \item $(d_{3})$ : $4x+y-6=0$ et $(d_{4})$ : $y=2x$.
\end{enumerate}

\exo[Deux droites à déterminer]
On considère les points $A(-1\,;1)$, $B(3\,;3)$ et $C(3\,;-2)$. La droite $(d)$ est la
droite $(AB)$ ; la droite $(d')$ passe par $C$ et a pour vecteur directeur
$\vec{u}\coord{-1}{2}$.
\begin{enumerate}
  \item Démontrer que $(d)$ a pour équation cartésienne $x-2y+3=0$.
  \item Déterminer une équation cartésienne de $(d')$.
  \item Démontrer que $(d)$ et $(d')$ sont sécantes.
  \item Calculer les coordonnées de leur point d'intersection $K$.
  \item Tracer $(d)$ et $(d')$ dans un repère et contrôler le résultat.
\end{enumerate}

\exo[Deux médianes d'un triangle]
On considère les points $A(2\,;6)$, $B(-4\,;0)$ et $C(8\,;0)$.
\begin{enumerate}
  \item Calculer les coordonnées des milieux $I$ de $[AB]$ et $J$ de $[AC]$.
  \item Démontrer que la droite $(CI)$ a pour équation $x+3y-8=0$.
  \item Déterminer une équation cartésienne de la droite $(BJ)$.
  \item Calculer les coordonnées du point d'intersection $G$ de $(CI)$ et $(BJ)$.
  \item Soit $M$ le milieu de $[BC]$. Démontrer que $G$ est sur la droite $(AM)$.
        Que représente le point $G$ pour le triangle $ABC$ ?
\end{enumerate}

\partiefiche{Équation réduite d'une droite par un système}

\exo[Par deux points]
Pour chaque couple de points, trouver l'équation réduite de la droite $(AB)$ à l'aide
d'un système.
\begin{enumerate}[label=\alph*)]
  \item $A(1\,;4)$ et $B(3\,;10)$.
  \item $A(-2\,;7)$ et $B(2\,;-1)$.
  \item $A(-1\,;0)$ et $B(2\,;1)$.
  \item $A(2\,;-1)$ et $B(5\,;1)$.
\end{enumerate}

\exo[Deux méthodes, un résultat]
On considère les points $A(-3\,;5)$, $B(1\,;-3)$ et $C(4\,;-9)$.
\begin{enumerate}
  \item Trouver l'équation réduite de la droite $(AB)$ à l'aide d'un système.
  \item Retrouver sa pente $m$ avec la formule $m=\dfrac{y_{B}-y_{A}}{x_{B}-x_{A}}$.
  \item Les points $A$, $B$ et $C$ sont-ils alignés ? Justifier.
\end{enumerate}

\exo[Salle de sport]
Une salle de sport fait payer un abonnement annuel, puis un prix fixe par séance. Avec
$10$ séances dans l'année, on paie $95$~€ en tout ; avec $25$ séances, $170$~€. Le prix
$y$, en euros, pour $x$ séances est donné par $y=mx+p$.
\begin{enumerate}
  \item Traduire l'énoncé par un système d'inconnues $m$ et $p$.
  \item Résoudre ce système.
  \item Que représentent $m$ et $p$ pour la salle de sport ?
  \item Combien paie-t-on pour $40$ séances ?
\end{enumerate}

\partiefiche{Aucune ou une infinité de solutions}

\exo[Combien de solutions ?]
Résoudre chaque système, puis dire par une phrase comment sont placées les deux droites
associées.
\begin{center}
  a)~${\sys{x-2y &= 3\\ -3x+6y &= 2}}$ \qquad
  b)~${\sys{2x+y &= 4\\ 6x+3y &= 12}}$\\[4pt]
  c)~${\sys{x+y &= 5\\ 2x+y &= 7}}$ \qquad
  d)~${\sys{4x-2y &= 6\\ -6x+3y &= -9}}$
\end{center}

\exo[Un coefficient inconnu]
Soit $a$ un nombre réel. On considère les droites $(d)$ : $2x+y-3=0$ et
$(d')$ : $ax+3y-1=0$.
\begin{enumerate}
  \item Donner un vecteur directeur de chaque droite.
  \item Exprimer le déterminant de ces deux vecteurs en fonction de $a$.
  \item Pour quelle valeur de $a$ les droites sont-elles parallèles ?
  \item Pour cette valeur de $a$, résoudre le système formé par les équations de
        $(d)$ et $(d')$. Conclure.
  \item Pour $a=4$, calculer les coordonnées du point d'intersection des droites.
\end{enumerate}

\exo[Le glacier et son perroquet]
Un glacier vend des glaces à la vanille à $2$~€ et des glaces au chocolat à $3$~€. On
note $x$ le nombre de glaces à la vanille vendues dans la journée et $y$ le nombre de
glaces au chocolat.
\begin{enumerate}
  \item Le soir, il dit à son perroquet : \og J'ai fait $150$~€ de recette. Si j'avais
        vendu toutes mes glaces deux fois plus cher, j'aurais fait $310$~€. \fg{}
        \og Impossible ! \fg{} répond le perroquet. A-t-il raison ? Justifier avec un
        système.
  \item Le glacier corrige : \og J'aurais fait $300$~€. \fg{} Le perroquet peut-il
        trouver $x$ et $y$ ? Justifier.
  \item Le lendemain, la recette est encore de $150$~€ ; avec les prix échangés
        ($3$~€ la vanille, $2$~€ le chocolat), elle aurait été de $160$~€. Combien de
        glaces de chaque parfum ont été vendues ?
\end{enumerate}

\end{multicols}

\end{document}
